\section{Contributions}
\label{sec:intro:contributions}
This thesis reduces the input rather than the model. It makes six contributions. (i) A
corpus of AIGs tiered by \emph{optimization state} (whether a target synthesis algorithm has
already been applied), extending OpenABC-D \citep{chowdhury2021openabc}
(\ref{sec:method:data}). (ii) A systematic benchmark of partitioning, sparsification and
summarization as a preprocessing step for GNN training on AIGs, measured on shrinkage, peak
memory, speed and predictive retention. (iii) AIG-specific adaptations of each family, built
from topological level, reconvergence structure and gate role, compared against their generic
counterparts at equal compression. (iv) A coarsening that is provably lossless for this
encoder \citep{bollen2023exact}, fixing a point of zero accuracy cost against which every
lossy method is measured (\ref{sec:method:architecture:exact}), together with the
architectural cost it imposes: normalisation and, where designs run to many levels,
node-level positional information (\ref{sec:method:summary:level}). (v) An evaluation of
cross-state inference, in which one set of weights is trained on reduced graphs and queried
on full ones. (vi) A controlled comparison of three train/test protocols, random,
recipe-disjoint and design-disjoint, ordered by the information each admits into the test set
(\ref{sec:method:experiment:splitting}). The hypothesis under test is that the design-disjoint
split alone measures structural inference, and that the excess accuracy of the other two is
recognition of circuits already seen in training.

The corpus is the enabling contribution: it draws on 55 source designs against OpenABC-D's 29,
and where OpenABC-D varies only the script applied to one starting circuit, it also varies the
optimization state that script is applied to. RQ3 and RQ5 depend on that distinction
(\ref{sec:relwork:ml4eda:datasets}).

